\section{Comparison with Causal Transformers}
\label{sec:comparison}

% This section is the analogue of Wu & Tu Section 3 ("Comparison with Transformers"), which is
% the LARGEST section of their paper -- about 2.4 pages, as long as their experiments. It had
% no counterpart in the first outline of this paper, and that was the biggest structural gap:
% in this format the conceptual comparison carries as much weight as the numbers, and it is
% also where the Looped baseline of Section 4 gets motivated rather than asserted.

\citet{wu2023probabilistic} show that MFVI in the encoder model reproduces the computation graph
of a transformer encoder, with the posterior over dependency heads playing the role of
self-attention. We ask the same question of the causal model, and the answer is the reason the
experiments in \Cref{sec:experiments} take the shape they do.

\subsection{Tensor form}
\label{sec:tensor-form}

\TODO{re-derive the update equations of \Cref{sec:inference} in tensor form, following the
structure of Wu \& Tu \S3.1 so the two papers can be read side by side. Their $Q_z^{(t)}$,
$Q_{h,c}^{(t)}$, $\mathcal{F}$, $\mathcal{G}$ notation should be reused verbatim wherever the
quantity is the same one --- deviating from it without reason costs the reader the comparison.}

\subsection{Single-channel update vs.\ masked self-attention}
\label{sec:vs-attention}

\TODO{the causal analogue of their \S3.2. The claim to establish: the single-channel update
under the causal chain reduces to scaled dot-product attention with a causal mask, and the
correspondence $Q_c \leftrightarrow W^Q$, $K_c = V_c \leftrightarrow W^K = W^V$ carries over.
State precisely where it stops holding --- that is more interesting than where it does.}

\subsection{The readout vs.\ a language-model head}
\label{sec:vs-lm-head}

\TODO{no counterpart in Wu \& Tu, because the encoder has no LM head. Show that
\Cref{eq:exact-readout} is a mixture of $\nchan$ softmaxes over a $\nlab$-dimensional space,
compare with the ordinary $\R^{\nlab\times|\Vocab|}$ head, and connect to
\citet{yang2018breaking}: the rank-$\nlab$ bottleneck is real, but a mixture of softmaxes is not
rank-limited in the way a single softmax is. Whether that helps is measured in
\Cref{sec:exp3}.}

\subsection{Full model comparison}
\label{sec:full-comparison}

\begin{figure*}[t]
\centering
\TODO{Figure 2, full width. Three computation graphs side by side --- (a) GPT-style decoder,
(b) Looped decoder, (c) causal probabilistic transformer --- in the visual style of Wu \& Tu's
Figure 3. This figure is what makes \Cref{tab:three-point} unnecessary to argue for in words.}
\caption{Computation graphs for a causal transformer, a Looped transformer, and a causal
probabilistic transformer.}
\label{fig:computation-graphs}
\end{figure*}

\TODO{prose. The differences to draw out, following their \S3.4:}

\begin{itemize}
  \item \textbf{No feed-forward structure.} Inherited from the encoder model, and the reason
        the global variables of \Cref{sec:variants} are worth measuring.
  \item \textbf{No residual connections or layer norm.} The unary scores are re-added at each
        iteration instead.
  \item \textbf{Parameters shared across iterations.} This is the property the Looped
        transformer \citep{dehghani2019universal, giannou2023looped} also has, and it is why a
        comparison against a standard decoder alone cannot attribute any difference to
        structure.
  \item \textbf{The output is a posterior, not a projection.} No parameter exists outside the
        factor graph (\Cref{sec:output}).
\end{itemize}

\begin{table}[t]
\centering
\small
\begin{tabular}{lcc}
\toprule
Model & \makecell{Weight\\sharing} & \makecell{Syntactic\\structure} \\
\midrule
Causal transformer & no  & no  \\
Looped transformer & yes & no  \\
Causal PT          & yes & yes \\
\bottomrule
\end{tabular}
\caption{The three models of \Cref{sec:experiments}. Looped vs.\ transformer isolates weight
sharing; causal PT vs.\ Looped isolates structure, which is the comparison this paper is about.}
\label{tab:three-point}
\end{table}

\Cref{tab:three-point} is what the experimental design follows from: the causal PT differs from
a standard decoder in two ways at once, so a two-model comparison cannot attribute a difference
to either. \NOTE{putting this table in Section 3 rather than Section 4 is deliberate --- it
makes the Looped baseline a consequence of the model comparison instead of an experimental
choice that a reviewer could read as cherry-picked.}
